%% kapa: exam=1 round=26 year=2015 period=2 form=A subject=accounting subject_ko=회계학 seq=23 number=63
%% source: Q-Net 감정평가사 기출문제 (cst00309) artlSeq=5138001 file=extracted/1008/5138001/제26회_감정평가사 1차_기출/제26회_감정평가사 1차_기출/2교시_감평_A형.pdf
%% conversion: kordoc 4.24.51 -> tools/gen1.py (auto); review=proofread-sonnet; flags=none
\begin{question}{63}
재무제표 표시의 일반사항에 관한 설명으로 옳지 않은 것은?
\begin{choices}
  \item 당해 기간의 기타포괄손익금액을 기능별로 분류해야 하며, 다른 한국채택국제회계기준서에 따라 후속적으로 당기손익으로 재분류되지 않는 항목과 재분류되는 항목을 각각 집단으로 묶어 표시한다.
  \item 경영진이 기업을 청산하거나 경영활동을 중단할 의도를 가지고 있지 않거나, 청산 또는 경영활동의 중단 외에 다른 현실적 대안이 없는 경우가 아니면 계속기업을 전제로 재무제표를 작성한다.
  \item 재무제표는 기업의 재무상태, 재무성과 및 현금흐름을 공정하게 표시해야 하며, 이를 위해서 ‘개념 체계’에서 정한 자산, 부채, 수익 및 비용에 대한 정의와 인식요건에 따라 거래, 그 밖의 사건과 상황의 효과를 충실하게 표현해야 한다.
  \item 사업내용의 유의적인 변화나 재무제표를 검토한 결과 다른 표시나 분류방법이 더 적절한 것이 명백한 경우, 한국채택국제회계기준에서 표시방법의 변경을 요구하는 경우가 아니면 재무제표항목의 표시와 분류는 매기 동일하여야 한다.
  \item 전체 재무제표는 적어도 1년마다 작성하되, 보고기간종료일의 변경으로 보고기간이 1년을 초과하거나 미달하게 되는 경우는 그렇지 않다.
\end{choices}
\end{question}
\answer{1}
